\documentclass[11pt]{article}
\usepackage[a4paper,margin=1in]{geometry}
\usepackage{amsmath,amssymb,amsthm}
\usepackage{hyperref}
\newtheorem{theorem}{Theorem}
\newtheorem{proposition}{Proposition}
\newtheorem{corollary}{Corollary}
\newtheorem{definition}{Definition}
\title{GST Causal-Graph Double Folds: Emergent Cohomological Dimension,
Hodge Charge, and Weight Defect}
\author{Research derivation (self-review; not Lean-kernel verified)}
\date{22 October 2026}
\begin{document}
\maketitle

\begin{abstract}
A dimension-free GST world can admit geometric realization maps without
assigning a dimension to the ambient causal ontology.  We make the user's
intuition testable using genuine projective nodal curves obtained by causal
identification of branches.  The dual graph of the identifications controls
a new weight-zero cohomology sector, and two independent causal loops create
\emph{rational diagonal Hodge-charge zero classes in even cohomological degree}
whose weight is strictly lower than their degree.  This persists under
arbitrarily many algebraic Tate dimension folds and survives stabilization
by every original smooth projective variety.  It is a theorem about the
mixed Hodge theory of singular realized worlds, \emph{not} a counterexample
to the rational Hodge conjecture for smooth projective varieties.
\end{abstract}

\section{Exact GST-to-geometry types and two diagnostic invariants}
The ontological source is the typed \texttt{GeneralSpace} of
\texttt{GSTGeneralSpace.lean}: points and associative paths, without
any built-in dimension or metric. The world-birth and event source is
\texttt{GSTGraphV2SelfExpandingCosmology.lean}; higher coherence cells in
\texttt{GSTGraphV2UnboundedHigherCausality.lean} have globular level $n$,
which is \emph{not} the complex dimension of a variety. The present
causal edge is realized by a specified algebraic gluing of branches.

For a graded piece of a mixed Hodge structure in cohomological degree
$k$, weight $w=p+q$, and Hodge bidegree $(p,q)$, define the
\emph{two-coordinate causal fold readout}
\[
  \chi=p-q,\qquad \delta=w-k=p+q-k.
\]
The signed GST/Hodge charge $\chi$ diagnoses diagonality.
The \emph{purity gap} $\delta$ diagnoses a distinct failure:
smooth projective complex varieties have $\delta=0$ in every
cohomological degree. Being diagonal $\chi=0$ without the correct
weight $\delta=0$ is insufficient even to instantiate the ordinary
degree-$2p$ rational Hodge-class target.

\begin{proposition}[Tate fold preserves both readouts]
An algebraic Tate shift by $j\geq0$ changes
\[
(k,p,q,w)\longmapsto(k+2j,p+j,q+j,w+2j),
\]
so $\chi$ and $\delta$ are both invariant under any finite succession
of such dimension folds.
\end{proposition}
\begin{proof}
Subtract the transformed coordinates:
$(p+j)-(q+j)=p-q$ and $(w+2j)-(k+2j)=w-k$.
Composition adds the shift indices.
\end{proof}

\section{General causal-graph source birth, not a label change}
Let $\Gamma$ be a finite connected undirected multigraph with
$r\geq1$ vertices, $e$ edges (loops and parallel edges allowed), and
\[
                  \beta=e-r+1\geq0.
\]
For each vertex take an independent copy of $\mathbb P^1_{\mathbb C}$.
For each edge choose two distinct marked smooth points, one on each
incident component (on the same component if a loop), and identify the
two points transversely to form an ordinary node.  The resulting
connected reduced projective nodal curve $C_\Gamma$ has normalization
\(\widetilde C_\Gamma=\coprod_{i=1}^r\mathbb P^1\).
The graph is the \emph{dual incidence graph of actual geometric
identifications}, not an arbitrary linear operator on Hodge labels.

\begin{theorem}[Causal graph cycles are emergent weight-zero holes]
For such $C_\Gamma$,
\[
 H^0(C_\Gamma,\mathbb Q)=\mathbb Q(0),\qquad
 H^1(C_\Gamma,\mathbb Q)\cong H^1(\Gamma,\mathbb Q)
       \cong\mathbb Q(0)^\beta,\qquad
 H^2(C_\Gamma,\mathbb Q)\cong\mathbb Q(-1)^r.
\]
The displayed $H^1$ is a pure weight-zero mixed Hodge summand in
cohomological degree one. It is born exactly when the causal graph has
a nontrivial cycle, not when it is a tree.
\end{theorem}
\begin{proof}
For the normalization $\nu:\widetilde C_\Gamma\to C_\Gamma$,
the difference of the two branch values at each node gives the exact
constant-sheaf sequence
\[
0\to\mathbb Q_{C_\Gamma}\to
\nu_*\mathbb Q_{\widetilde C_\Gamma}\to
\bigoplus_{\mathrm{node}\in E(\Gamma)}\mathbb Q_{\mathrm{node}}\to0.
\]
Since the normalized components are rational,
$H^1(\widetilde C_\Gamma,\mathbb Q)=0$.  The long exact sequence
contains
\[
0\to\mathbb Q\xrightarrow{\Delta}\mathbb Q^r
\xrightarrow{B_\Gamma}\mathbb Q^e
\to H^1(C_\Gamma,\mathbb Q)\to0,
\]
where $B_\Gamma$ is the oriented edge--vertex incidence matrix.
Connectedness makes $\operatorname{rank} B_\Gamma=r-1$ (a
spanning-tree argument, valid with loops), so the cokernel has rank
$e-r+1=\beta$.  This is also the graph cohomology $H^1(\Gamma,\mathbb Q)$.
The normalization sequence is compatible with mixed Hodge structures:
its graph quotient comes from weight-zero $H^0$ of points and
components, hence has weight zero and type $(0,0)$.
In degree two normalization gives $\mathbb Q(-1)^r$.
\end{proof}

\section{All-dimensional causal fold laws}
For every integer $N\geq0$ put
\[
       Y_{\Gamma,N}=C_\Gamma\times\mathbb P^N.
\]
This is an \emph{actual} projective complex realized world with complex
dimension $N+1$; it is singular when $\Gamma$ has any nodes.
Write $\mathcal H(t;u,v,s)$ for the degree--Hodge--weight polynomial,
where $t$ counts ordinary cohomological degree, $u^pv^q$ a Hodge
bidegree on $\operatorname{Gr}^W_{p+q}$, and $s$ the weight.
With $a=t^2s^2uv$, mixed-Hodge Kunneth yields the closed all-$N$ law
\[
\boxed{\quad
  \mathcal H(Y_{\Gamma,N};t,u,v,s)
  =(1+\beta t+r\,a)\sum_{j=0}^N a^j.
 \quad} \tag{1}
\]
In particular, at each $0\leq j\leq N$,
\[
\boxed{\ \operatorname{Gr}^{W}_{2j}
 H^{2j+1}(Y_{\Gamma,N},\mathbb Q)
 \cong\mathbb Q(-j)^\beta,\quad
 (p,q)=(j,j),\quad(\chi,\delta)=(0,-1).\ }\tag{2}
\]
Taking $t=-1,s=1$ gives independently the compact projective
Hodge--Deligne polynomial
$E(Y_{\Gamma,N};u,v)=(1-\beta+r\,uv)
\sum_{j=0}^N(uv)^j$.

\begin{theorem}[Two independent causal folds produce even-degree
charge-zero, non-pure classes]
Let $\Gamma_1,\Gamma_2$ be two connected finite multigraphs with
cycle ranks $\beta_1,\beta_2>0$, and with all normalization components
$\mathbb P^1$. Define
$Z_N=C_{\Gamma_1}\times C_{\Gamma_2}\times\mathbb P^N$.
For every $0\leq j\leq N$,
\[
\boxed{\quad
\operatorname{Gr}^{W}_{2j} H^{2j+2}(Z_N,\mathbb Q)
\cong
 H^1(\Gamma_1,\mathbb Q)\otimes
 H^1(\Gamma_2,\mathbb Q)(-j)
\cong\mathbb Q(-j)^{\beta_1\beta_2}.
\quad}\tag{3}
\]
Every nonzero class in this piece has
\[
    k=2j+2,\quad (p,q)=(j,j),\quad
    w=2j,\quad\chi=0,\quad\delta=-2.
\]
Thus the causal fold has created a genuinely rational, diagonal,
even-degree cohomological sector which the \emph{charge-only}
test accepts but the purity test excludes from the usual
smooth-projective Hodge conjecture.
\end{theorem}
\begin{proof}
The rational mixed-Hodge Kunneth decomposition contains
$H^1(C_{\Gamma_1})\otimes H^1(C_{\Gamma_2})
\otimes H^{2j}(\mathbb P^N)$ in degree $2j+2$.
Its weight is $0+0+2j=2j$, its type is $(j,j)$, and its rank is
$\beta_1\beta_2$. All other degree-$2j+2$ summands of
$C_{\Gamma_1}\times C_{\Gamma_2}\times\mathbb P^N$
have weight $2j+1$ or $2j+2$: each normalized rational-curve
factor contributes either $H^0$ of degree/weight $(0,0)$ or
$H^2$ of degree/weight $(2,2)$, while a graph $H^1$ contributes
$(1,0)$. Weight deficit $-2$ therefore occurs precisely in the
two-graph $H^1\otimes H^1$ term. This establishes the equality
of the weight graded piece, not merely a lower bound.
\end{proof}

\begin{corollary}[The phenomenon survives every arbitrary original $X$]
For each nonempty connected smooth projective complex variety $X$,
$X\times Z_N$ contains, in degree $2j+2$, the nonzero rational
weight-$2j$ direct summand of (3), embedded through
$H^0(X,\mathbb Q)=\mathbb Q$. Thus the causal-fold diagnostic works
with arbitrary original $X$ and arbitrarily high additional emergent
geometric dimension \emph{in a new singular world}. It neither
modifies nor refutes algebraicity of a prescribed class on the
untouched $X$.
\end{corollary}
\begin{proof}
Apply mixed-Hodge Kunneth, taking the $H^0(X)$ summand.
\end{proof}

\section{Independent charge leakage and exact interpretation}
The two-channel invariant $(\chi,\delta)$ distinguishes:
(i) holomorphic algebraic Tate folds: $(\chi,\delta)$ is preserved;
(ii) topological nonholomorphic operators on fixed smooth
manifolds: $\chi$ may leak while $\delta$ remains zero because
cohomology stays pure in each fixed geometric realization;
(iii) actual singular causal graph folds: the graph cycle creates
$\delta=-1$, and two independent folds create
$(\chi,\delta)=(0,-2)$ in even degree.
The sign of $\delta$ is a genuine weight-vs-degree witness, not
the label of an abstract GST cell.

A tree has $\beta=0$ and contributes no new weight-zero graph
cohomology, although it may still produce singular nodes.
One causal cycle gives one new independent hole at all Tate levels.
Two independent graph cycles multiply the rank rather than merely
adding it: $\beta_1\beta_2$.
This produces a dimension/causality equation testable on the
\emph{realized space itself}, unlike redefining an abstract Hodge map.

\section{Exact Wolfram finite checks; scope and unresolved bridge}
A stateless Wolfram evaluation tested the edge--vertex incidence
ranks of a 3-vertex tree $(e=2,\beta=0)$, a triangle
$(e=3,\beta=1)$, a bouquet of three loops
$(r=1,e=3,\beta=3)$ and a 4-vertex 5-edge graph
$(e=5,\beta=2)$.  The computed matrix ranks were
$2,2,0,3$, as the spanning-tree derivation predicts.

The degree/weight/Hodge generating-function coefficients (1)--(2)
were checked in 84 parameter pairs for $1\leq\beta\leq4$,
$0\leq N\leq5$ and every $0\leq j\leq N$.
For two independent rational-curve node families, equation (3)
was checked in 324 parameter tuples
$1\leq\beta_1,\beta_2\leq3$, $0\leq N\leq7$,
$0\leq j\leq N$.
All checks returned exact symbolic equality in Wolfram Language.
These checks are consistency tests; the proofs above are uniform
and do not rely on the finite bound. No new Lean code or GitHub
Actions compiler run is claimed for the results of this paper.

\paragraph{Critical nonclaim.}
A singular projective world has a mixed Hodge structure, not the
pure degree-$k$ Hodge structure assumed in the classical rational
Hodge conjecture. A charge-zero class of weight $2j$ occurring in
cohomological degree $2j+2$ \emph{does not become} a counterexample
to rational algebraicity on smooth projective $X$.
For an unconditional original-$X$ proof the independent missing
map is a genuine cycle
$Z_\alpha\in CH^p(X)_{\mathbb Q}$ with
$\mathrm{cl}^{\mathrm{can}}_X(Z_\alpha)=\alpha$.
The graph-fold weight law neither assumes nor constructs this map.
It instead supplies a new, precise causal birth/dimensional purity
discriminator for future GST realized operations.

\end{document}
